\subsection{相对二次体积形式}\label{sec:15c}
令 $\pi:E\to B$ 为带奇点的纤维化。相对二次体积形式丛为
\[
\mathcal K_{E/B}^2=\pi^*\lambda_{\CC}^{\mathrm{top}}(TB)^{\otimes2}\otimes
\lambda_{\CC}^{\mathrm{top}}(TE)^{\otimes-2}.
\]
假设有该丛的处处非零截面 $\eta_{E/B}^2$；最简单的取得方式，是用总空间和底空间上的二次体积形式之商 $\eta_E^2/\eta_B^2$，但一般情形不必如此。若 $b$ 是 $\pi$ 的正则值，则短正合列 $0\to T(E_b)\to TE|E_b\to\pi^*TB_b\to0$ 诱导典范同构
\[
\mathcal K_{E/B}^2|E_b\cong\mathcal K_{E_b}^2.
\]
因此，$\eta_{E/B}^2$ 在光滑纤维上确定二次复体积形式 $\eta_{E_b}^2$。类似地，得到相对平方相位映射 $\alpha_{E/B}:\operatorname{Gr}(TE^v)\to S^1$，在临界点之外定义，其在每个光滑纤维上的限制与相应的 $\alpha_{E_b}$ 一致。在此情况下，平行移动映射\eref{15.2}取得自然的分次，即函数 $\gamma_{h_\beta}^\#:\operatorname{Gr}(E_{\beta(0)})\to\R$，使
\begin{equation}\label{eq:15.8}
\exp(2\pi i\,\gamma_{h_\beta}^\#(\Lambda))
=\frac{\alpha_{E_{\beta(1)}}(Dh_\beta(\Lambda))}{\alpha_{E_{\beta(0)}}(\Lambda)}.
\end{equation}
从 $\Id_{E_{\beta(0)}}$ 的零函数出发，对所有 $s$ 连续延伸到 $\beta|[0,s]$，便得到这些函数。在 $\beta$ 为闭路径的特殊情形，结果是 $h_\beta$ 成为\secref{12i}所定义的分次辛自同构。\znote{216}{相对二次体积形式的作用是为纤维化提供“沿纤维的典范数据”。因为纤维的切空间由正合列 $0\to T(E_b)\to TE|E_b\to\pi^*TB_b\to0$ 给出，$\mathcal K_{E/B}^2$ 限制到每条光滑纤维上就是 $\mathcal K_{E_b}^2$；所以一个相对截面同时为所有光滑纤维指定二次复体积形式，不必逐条纤维选择。相位 $\alpha_{E/B}$ 是 $\mathcal K_{E/B}$ 的平方根在 Lagrangian Grassmann 流形上的相位，它使平行移动能够带上分次：\eref{15.8}把端点相位之比的对数取为 $\gamma^\#_{h_\beta}$，沿路径连续延拓即得，而闭路径情形下所得正是分次辛自同构。若没有相对数据，就只能逐条纤维比较，无法保证与平行移动相容。[来源：原书 §(15c)，印刷页 216。]}

\subsection{Lefschetz 纤维化}\label{sec:15d}
令 $S$ 为具有非空边界的连通紧 Riemann 曲面。选择一形式 $\theta_S$，满足 $\theta_S|\partial S>0$ 及 $\omega_S=d\theta_S>0$，就使其成为带角正合辛流形。\emph{（正合）Lefschetz 纤维化}是带奇点的纤维化 $\pi:E^{2n+2}\to S$，其临界点一般，也称为非退化，且局部可积。这意味着 $I_E$ 在 $\operatorname{Crit}(\pi)$ 的某个邻域中可积，且 $D\pi$ 作为复线性映射丛 $\Hom_{\CC}(TE,\pi^*TS)$ 的截面与零截面横截。第二个条件等价于每个临界点处的复 Hessian $D^2\pi$ 非退化。此外，还假设每个纤维中至多有一个临界点，故投影 $\operatorname{Crit}(\pi)\to\operatorname{Critv}(\pi)$ 为双射；最后这一假设只是为方便起见，很容易去除，代价是后面的记号稍显繁琐。

非退化性蕴含临界点孤立，故 $\operatorname{Crit}(\pi)$ 是 $\operatorname{int}(E)$ 的有限子集，类似地，$\operatorname{Critv}(\pi)$ 是 $\operatorname{int}(S)$ 的有限子集。在每个
\SourcePage{225}
临界点及其临界值附近，存在局部全纯坐标，使 $\pi$ 成为标准二次映射
\begin{equation}\label{eq:15.9}
Q(x)=x_1^2+\cdots+x_{n+1}^2.
\end{equation}
一般而言，$\omega_E$ 在这些坐标中并不标准。不过，可以找到沿 $\partial^hE$ 行为良好的纤维化变形，事实上它局限在临界点附近，使在变形的另一端，Kähler 形式在给定的全纯 Morse 坐标中成为标准形式。这使下述基本模型更为重要：

\StatementNumber{4}\begin{example}\label{st:15.4}
令 $Q:\CC^{n+1}\to\CC$ 如\eref{15.9}，并令 $k:\CC^{n+1}\to\R^+$ 为函数 $k(x)=(|x|^4-|Q(x)|^2)/4$。对固定的 $r,s>0$，定义
\begin{equation}\label{eq:15.10}
E=\{x\in\CC^{n+1}:|Q(x)|\leq r,\ k(x)\leq s\},
\end{equation}
赋予 $\CC^{n+1}$ 上标准辛形式的限制、它的标准原始形式 $(i/4)(z\,d\bar z-\bar z\,dz)$、给定复结构 $I_E=i$，以及限制 $Q$ 所得的映射 $\pi:E\to S=rD$。边界面为
\[
\partial^vE=\{x\in E:|\pi(x)|=r\},\qquad
\partial^hE=\{x\in E:k(x)=s\}.
\]
截断函数 $k$ 的选取使 $TE^h$ 平行于 $\partial^hE$。为验证这一点，注意 $TE_x^h$ 在 $\CC$ 上由 $(\nabla Q)_x=2\bar x$ 生成，再核对 $dk_x(\bar x)=0$、$dk_x(i\bar x)=0$ 即可。每个非奇异纤维 $E_z$，$z\ne0$，都与余切向量在标准度量下长度 $\leq\sqrt{s}$ 的子集 $B_{\sqrt{s}}^*S^n\subset T^*S^n$ 辛同构。明确地，
\[
B_{\sqrt{s}}^*S^n=\{(u,v)\in\R^{n+1}\times S^n:
\langle u,v\rangle=0,\ |u|\leq\sqrt{s}\},
\]
带辛形式 $du\wedge dv$；当 $z>0$ 时，同构 $E_z\to B_{\sqrt{s}}^*S^n$ 由下式给出：
\[
\phi_z(x)=(-\operatorname{im}(x)|\operatorname{re}(x)|,
\operatorname{re}(x)|\operatorname{re}(x)|^{-1}).
\]
遗憾的是，尽管负 Liouville 向量场，即负径向向量场，确实沿 $\partial E$ 指向内部，$E$ 却并非弱 $I_E$ 凸；纤维是弱凸的，总空间则不是。因此，这还不完全是这里定义的正合 Lefschetz 纤维化的例子，虽然从纯辛的观点看，它具有全部所需性质。当然，可以改变 $I_E$ 来改善这种情况，但实际上没有必要，因为最终 $E$ 只用作局部模型。
\end{example}

\StatementNumber{5}\begin{note}\label{st:15.5}
辛 Lefschetz 纤维化的其他定义，例如见\bcite{39}或\bcite{64}。通常假设底空间和纤维均闭，这在技术上略为简单。局部标准形\eref{15.9}的存在是全纯 Morse 引理\bcite{8}的推论；更准确地说，对底空间上任意选取的全纯坐标，都可在总空间上找到坐标，使 $\pi=Q$。使这些坐标中的辛结构标准化的变形，构造于\bcite{131}引理 1.6。\stref{例题}{15.4}的几何在\bcite{131}第 1.2 节中有详尽讨论。
\end{note}
